% !TeX root = ../教材题库_全部条目_转写工作稿.tex
% Source: XJ-B1-C02-RV-P063-Y1035
% 第2章 / 章末复习 复习题二 / 教材 p.63 / PDF 第70页
% OCR未校对：公式与文字均可能有误；未完成数学TeX转写；禁止直接发给学生
% 原题及相邻原解截图（核对用）：../题目原图/XJ-B1-C02-RV-P063-Y1035-01.webp
% 本文件只在转写完成后作为题干使用；OCR 条目尚不是可靠题干。
\begin{bankitem}{复习题二 第17题}{第2章 / 章末复习 复习题二 / 教材 p.63 / PDF 第70页}{XJ-B1-C02-RV-P063-Y1035}{OCR未校对：公式与文字均可能有误；未完成数学TeX转写；禁止直接发给学生}
\begin{ocrdraft}
[教材p.63 / PDF第70页]
\par
17．
\par
为了在夏季降温和冬季供暖时减少能源损耗，房屋的屋顶和外墙需要建
\par
造隔热层．某幢建筑物要建造可使用20年的隔热层，每厘米厚的隔热层建造成本
\par
为6万元．该建筑物每年的能源消耗费用C（万元）与隔热层厚度工（cm）满足如下
\par
关系式：
\par
（0工10）．
\par
3工十5
\par
若不建隔热层，每年能源消耗费用为8万元．设万元）为隔热层建造费用
\par
与20年的能源消耗费用之和．
\par
（D求差的值及3，的表达式．
\par
（2）隔热层修建多厚时，总费用3，最小？并求最小值．
\end{ocrdraft}
\end{bankitem}
